\section{System Architecture}
\label{sec:architecture}

\subsection{Operational Setting and Requirements}
\label{sec:requirements}

The deployment environment is municipal rescue-service incident response in a country with substantial unpopulated forested terrain. Incidents are managed by a small set of vehicles, typically two to ten per incident, operating up to several hours away from the regional dispatch centre. Each vehicle is staffed by personnel using rugged tablets to view and update operational information. The incident is led by a designated command vehicle whose role is operationally fixed at dispatch and reassigned only by explicit operator decision. The publicly available rescue-service guidance summarised in Section~\ref{sec:introduction} establishes that incident-planning information must be available without continuous connectivity, that secrecy-protected and non-secret material must be handled with different operational discipline, and that incident plans are layered, ongoing operational artefacts rather than static documents~\cite{rsg-rad201, rsg-kommunalplan}.

We distinguish externally motivated requirements from design assumptions and deployment constraints. The distinction matters because an operational command authority does not logically require a single-writer storage design; aligning the storage authority boundary with the operational authority boundary is the design choice made in this paper.

\begin{itemize}
  \item \textbf{R1 (Offline operation).} Each tier must continue functioning on previously synchronised data without connectivity to any other tier.
  \item \textbf{R2 (Durable forward-only submission).} Field observations entered on tablets must never be lost, including across vehicle restart and prolonged partition.
  \item \textbf{R3 (Attributable command authority).} Decisions on resource assignments, status transitions, and hazard annotations must be attributable to a single operationally designated authority at every instant of the incident.
  \item \textbf{R4 (Predictable network behaviour).} The wireless mesh substrate must not cause cascading failures (broadcast storms, routing loops) when nodes join, leave, or partition.
  \item \textbf{R5 (No parallel identity).} Authentication, authorisation, and audit must reuse the rescue-service organisation's existing certificate authority and identity provider; no parallel identity stack may be introduced.
\end{itemize}

\begin{table}[t]
\centering
\caption{Traceability from adopted constraints to architectural consequences.}
\label{tab:req-trace}
\footnotesize
\begin{tabularx}{\columnwidth}{@{}l l X@{}}
\toprule
\textbf{Item} & \textbf{Type} & \textbf{Result} \\
\midrule
Offline operation & Operational requirement & Edge cache and local services. \\
Durable field submission & System requirement & Tablet/edge outbox and retry. \\
Attributable command authority & Operational assumption & Single command-side journal writer. \\
Predictable network behavior & Deployment constraint & Mesh as routed transit; no stretched user VLAN. \\
Existing identity integration & Integration constraint & mTLS/CA integration and mapped identities. \\
\bottomrule
\end{tabularx}
\end{table}

The five design principles in Section~\ref{sec:principles} are derived from this traceability rather than asserted independently. The mapping is explicit: \textit{offline-first} from R1; \textit{single incident writer} from R3 and the chosen storage-authority boundary; \textit{forward-only edits} from R2 directly and R3 indirectly; \textit{mesh as transit only} from R4 directly and R5 indirectly, because stretched VLANs would extend the trust surface across the mesh; \textit{existing identity} from R5.

\subsection{Design Principles}
\label{sec:principles}

The system is built around five principles:
\begin{itemize}
  \item \textbf{Offline-first} (R1): all previously synchronised data remains available without connectivity.
  \item \textbf{Single incident writer} (R3, R2): exactly one command node sequences and commits incident events.
  \item \textbf{Forward-only edits} (R2, R3): tablets and responder vehicles enqueue edits locally and forward them upstream.
  \item \textbf{Mesh as transit only} (R4, R5): the wireless mesh carries traffic between edge nodes, but does not expose user VLANs end to end.
  \item \textbf{Existing identity} (R5): the deployment integrates with the authority's established identity and certificate infrastructure.
\end{itemize}

\subsection{Tier 1: Regional Core}

The regional core hosts the authoritative master-data platform in the target deployment. It uses a spatial database, a publication pipeline that stages, validates, and promotes geodata into versioned offline packages, and a tile-serving stack for map data. A reverse proxy terminates TLS and routes API traffic to internal services. A VPN overlay connects enrolled vehicles when WAN is present. Published packages consist of map bundles, compressed attachment archives, and SHA-256 manifests that edge nodes can pull over HTTP range requests. Section~\ref{sec:implementation} separates which of these elements are implemented and which are deployment scaffolding.

\subsection{Tier 2: Vehicle Edge}

The vehicle edge tier is the architectural center of the paper. Each vehicle carries the same edge hardware profile, but nodes operate in one of two software roles:
\begin{itemize}
  \item \textbf{Command vehicle:} commits incident events, assigns sequence numbers, materializes incident state, and backhauls to core when WAN is present.
  \item \textbf{Responder vehicle:} caches master data, serves local tablets, and queues field edits in an outbox until they can be forwarded to command.
\end{itemize}

This role split is intentional. It preserves hardware interchangeability while keeping the incident journal under one authority-bearing node at any time. Any responder can be promoted to command by configuration and operational procedure.

\subsection{Tier 3: Field Tablets}

Field tablets are leaf clients attached to a local vehicle only. They use an encrypted local database and local tile cache, and they communicate over HTTPS with the edge node in their vehicle. Tablets have no direct mesh participation and no direct internet dependency. This sharply limits the exposed trust and connectivity surface.

\subsection{State Classes and Consistency Contracts}
\label{sec:state-taxonomy}

The single-writer constraint applies to authority-bearing incident-journal entries, not to every byte in the system. Figure~\ref{fig:state-contracts} separates the main state classes and the contract each one uses.

\begin{figure}[t]
\centering
% \begin{tikzpicture}[
%   x=1cm,y=1cm,
%   every node/.style={font=\scriptsize},
%   head/.style={draw, fill=gray!25, font=\scriptsize\bfseries, align=center, minimum height=0.42cm, inner sep=1pt},
%   state/.style={draw, fill=gray!10, align=left, text width=1.75cm, minimum height=0.48cm, inner sep=2pt},
%   src/.style={draw, align=center, text width=1.55cm, minimum height=0.48cm, inner sep=2pt},
%   contract/.style={draw, align=left, text width=4.25cm, minimum height=0.48cm, inner sep=2pt},
%   emph/.style={contract, double, double distance=0.5pt}
% ]
% \node[head, text width=1.75cm] at (0.90,0) {State};
% \node[head, text width=1.55cm] at (2.65,0) {Authority / source};
% \node[head, text width=4.25cm] at (5.60,0) {Synchronization contract};

% \node[state] at (0.90,-0.55) {Reference data};
% \node[src] at (2.65,-0.55) {regional core};
% \node[contract] at (5.60,-0.55) {one-way core-to-edge replication; eventual convergence};

% \node[state] at (0.90,-1.15) {Immutable artifacts};
% \node[src] at (2.65,-1.15) {publisher / core};
% \node[contract] at (5.60,-1.15) {manifest/package/hash validation where implemented};

% \node[state] at (0.90,-1.75) {Field observations};
% \node[src] at (2.65,-1.75) {tablet / responder};
% \node[contract] at (5.60,-1.75) {locally durable outbox; at-least-once forwarding; idempotent accept};

% \node[state] at (0.90,-2.35) {Command decisions};
% \node[src] at (2.65,-2.35) {command edge};
% \node[emph] at (5.60,-2.35) {single-writer applies only here; total order by event\_seq};

% \node[state] at (0.90,-2.95) {Materialized state};
% \node[src] at (2.65,-2.95) {journal};
% \node[contract] at (5.60,-2.95) {derived/rebuildable; not authoritative};

% \node[state] at (0.90,-3.55) {Audit events};
% \node[src] at (2.65,-3.55) {edge / core};
% \node[contract] at (5.60,-3.55) {eventual forwarding; validation incomplete};
% \end{tikzpicture}
\scriptsize
\setlength{\tabcolsep}{3pt}
\begin{tabularx}{\columnwidth}{@{}p{0.24\columnwidth}p{0.20\columnwidth}X@{}}
\toprule
\textbf{State} & \textbf{Authority/source} & \textbf{Synchronization contract} \\
\midrule
Reference data & regional core & One-way core-to-edge replication; eventual convergence. \\
Immutable artifacts & publisher / core & Manifest, package, or hash validation where implemented. \\
Field observations & tablet / responder & Locally durable outbox; at-least-once forwarding; idempotent accept. \\
Command decisions & command edge & Single-writer rule applies only here; total order by \textit{event\_seq}. \\
Materialized state & journal & Derived and rebuildable; not authoritative. \\
Audit events & edge / core & Eventual forwarding; validation incomplete. \\
\bottomrule
\end{tabularx}
\caption{State classes use different synchronization contracts. The single command writer applies to authority-bearing journal entries, not to every object in the system; materialized state is derived, and audit forwarding validation remains incomplete.}
\label{fig:state-contracts}
\end{figure}

\subsection{Network Segmentation}

Each vehicle is intended to carry a router enforcing a five-VLAN segmentation model keyed on a per-vehicle identifier. Tablets reach only the local \texttt{ops-api} endpoint over HTTPS, user VLANs are not extended across the mesh, and the mesh operates as Layer~3 transit for application-layer traffic between edge nodes. Figure~\ref{fig:topology-scope} separates this target deployment from the narrower Docker/Python path measured in Section~\ref{sec:validation}.

\begin{figure}[t]
\centering
\begin{tikzpicture}[
  node distance=4mm,
  every node/.style={font=\scriptsize},
  svc/.style={draw, rounded corners=1pt, align=center, minimum height=5mm, inner sep=2pt},
  core/.style={svc, fill=gray!15, text width=22mm},
  edge/.style={svc, text width=19mm},
  tablet/.style={svc, text width=13mm, minimum height=4mm},
  emu/.style={svc, fill=white, text width=16mm, minimum height=5mm},
  targetlink/.style={->, >=Stealth, dashed, semithick},
  measuredlink/.style={->, >=Stealth, semithick},
  boundary/.style={draw, rounded corners=2pt, inner sep=3pt}
]
\node[core] (core) at (3.75,1.0) {Regional\\core};
\node[edge] (cmd) at (2.0,-0.55) {Command\\K430};
\node[edge] (resp) at (5.45,-0.55) {Responder\\K430};
\node[tablet] (tab1) at (2.0,-1.85) {Tablet};
\node[tablet] (tab2) at (5.45,-1.85) {Tablet};

\draw[targetlink] (tab1) -- node[right, font=\tiny] {HTTPS/mTLS} (cmd);
\draw[targetlink] (tab2) -- node[right, font=\tiny] {HTTPS/mTLS} (resp);
\draw[targetlink] (cmd) -- node[above, font=\tiny] {Rajant transit} (resp);
\draw[targetlink] (cmd) -- node[left, font=\tiny] {WG/HTTPS} (core);
\draw[targetlink] (resp) -- node[right, font=\tiny] {WG when online} (core);

\node[boundary, fit=(core)(cmd)(resp)(tab1)(tab2), label={[font=\scriptsize]above left:target deployment}] (targetbox) {};

\node[emu, text width=14mm] (stub) at (0.45,-4.0) {Python\\stub};
\node[emu, text width=14mm] (ops) at (2.10,-4.0) {resp.\\ops-api};
\node[emu, text width=14mm] (accept) at (3.95,-4.0) {cmd.\\accept};
\node[emu, text width=14mm] (journal) at (5.55,-4.0) {cmd.\\journal};
\node[emu, text width=14mm] (syncapi) at (7.15,-4.0) {core\\sync-api};

\draw[measuredlink] (stub) -- node[above, font=\tiny] {POST event} (ops);
\draw[measuredlink] (ops) -- node[above, font=\tiny] {outbox batch} (accept);
\draw[measuredlink] (accept) -- node[above, font=\tiny] {seq} (journal);
\draw[measuredlink] (journal) -- node[above, font=\tiny] {backfill} (syncapi);

\node[boundary, fit=(stub)(ops)(accept)(journal)(syncapi), label={[font=\scriptsize]above left:reported emulation}] {};

\node[align=left, font=\tiny, anchor=west] at (0.05,-5.0) {solid = measured Docker/Python path; dashed = target path not measured};
\node[align=left, font=\tiny, anchor=west] at (0.05,-5.35) {tablet is a leaf client; mesh is transit, not tablet participation};
\end{tikzpicture}

\caption{Target deployment layers and the narrower path exercised in the reported emulation. The Docker/Python path evaluates the outbox, command journal, and core-forwarding behavior; physical mesh, VPN, mTLS, and Android client behavior are target-deployment elements not measured here.}
\label{fig:topology-scope}
\end{figure}
